%!TeX root=csp-blueprint.tex
\section{What a formalization can claim}\label{sec:bp-targets}

The dichotomy theorem, as stated in the proof paper, is conditional on a model
of computation: ``solvable in polynomial time'' quantifies over an encoding of
instances as strings, a machine, and a polynomial bound on its step count, and
``\textup{NP}-complete'' additionally quantifies over all problems in
\textup{NP} and over polynomial-time many-one reductions. Mathlib supplies
Turing machines (\texttt{Mathlib/Computability/TuringMachine/}) and a notion of
computability, but it has \emph{no} complexity class, no polynomial-time
predicate, and no notion of polynomial-time reduction. So a formalization must
say which of the statements it is proving, and the proof paper's five layers
are the answer.

\subsection{Where the work actually is}\label{sec:program-axis}

An earlier draft of this section concluded from that gap that the
complexity-theoretic half should be deferred because building it would dominate
the project. That inference was wrong, and it is worth recording why, because
the corrected version is what organises the rest of this document.

The complexity vocabulary itself --- languages, a cost model, $\mathrm P$,
$\mathrm{NP}$ by verifiers, $\le_{p}$, \textup{NP}-hardness,
\textup{NP}-completeness, and a generic \textup{NP}-hard problem --- has since
been built, in about $850$ lines. That is consistent with the two existing
Cook--Levin formalizations, whose class layers are $702$ lines (Isabelle) and
$992$ lines (Coq) against totals of $56\,514$ and $\approx 16\,500$. Defining
the classes is cheap in every development that has done it.

What is expensive is not complexity theory. It is any statement that requires
one to \emph{exhibit a program}. In a formalization, ``$f$ is computable in
polynomial time'' is not a property one observes of a mathematical function; it
is an existential over programs, discharged only by writing one and proving a
bound on its step count. The dividing line runs along that axis, and it cuts
across the algebra/complexity split rather than agreeing with it:

\begin{center}
\begin{tabular}{@{}lll@{}}
\hline
Target & Exhibit a program? & Cost \\
\hline
\textup{(T0)} the algebraic core & no & the mathematics \\
\textup{(T1)} $\textsc{Solve}$ is correct & no (but see Formalization note~\ref{haz:t1-computable}) & moderate \\
\textup{(T2)} $\textsc{Solve}$ is polynomial & \textbf{yes --- the largest one here} & \textbf{large} \\
\textup{(H0)} $\Gamma$ pp-interprets \textsc{Nae} & no & the mathematics \\
\textup{(H1)} pp-interpretation $\Rightarrow\le_{p}$ & yes, a syntactic substitution & small \\
\hline
\end{tabular}
\end{center}

The blueprint is written for \textup{(T0)}, \textup{(T1)} and \textup{(H0)}.
\textup{(H1)} is a small addition on top; \textup{(T2)} is neither small nor
mathematics, but it is not \emph{unrelated} to the dichotomy either --- it is
the assertion that the algorithm one has just proved correct is also fast, and
it is the second-largest component of the project.

\begin{hazard}[\textup{(T1)} is vacuous unless $\textsc{Solve}$ computes]
\label{haz:t1-computable}
``$\textsc{Solve}$ written as a total function'' hides a requirement that no
cost model would catch. If $\textsc{Solve}$ is defined using classical choice
--- for instance by deciding ``does $D_{x}$ have a nonempty proper binary
absorbing subuniverse?'' with \texttt{Classical.dec} --- then
$\textsc{Solve}(\mathcal I)=\texttt{true}\iff\mathcal I$ has a solution is a
true theorem that says nothing algorithmic, because $\textsc{Solve}$ does not
reduce. To mean what it appears to mean, \textup{(T1)} needs every predicate
the algorithm branches on to carry a genuine decision procedure.

That is a real obligation with a locatable cost. Absorption is defined by an
existential over \emph{terms} --- an infinite type --- so ``$C$ is a binary
absorbing subuniverse of $B$'' is not decidable on its face. It becomes
decidable through a bound on the arity of a sufficient witness, which is a
theorem and not a definition. For the types actually used here the bound is
cheap and should be taken that way rather than from a general absorption-arity
theorem: $\TBA$ has a binary witness by definition, and central absorption has
a ternary witness by \inproof{lem:central-ternary}. Only finitely many binary
and ternary operations then need be checked, together with finitely many
congruences and subsets. The same remark applies to the searches inside
$\textsc{SolveLinear}$.
\end{hazard}

\begin{hazard}[\textup{(T2)} is not a wrapper]\label{haz:t2-is-not-small}
Two things make \textup{(T2)} much larger than the phrase ``complexity
wrapper'' suggests.

\emph{The program is enormous.} \textup{(T2)} requires $\textsc{Solve}$ ---
\textsc{ForceConsistency}, the search for a strong subuniverse,
\textsc{SolveLinear}, and the mutual recursion between them --- to be expressed
in a cost model, with a proved polynomial bound. Every existing measurement
says this kind of work dominates: it is the $50\,000$ lines of Turing-machine
engineering in the Isabelle Cook--Levin development, against $702$ for the
classes themselves.

\emph{The statement is non-uniform.} Because $\Gamma$ is fixed, the polynomial
may depend on $\Gamma$, and the searches over subsets of $A$, over congruences
on $A$, and over term operations of bounded arity are all constant-time. But
``constant'' here means the search is compiled away into a program whose
\emph{size} depends on $|A|$. So $\textsc{Solve}$ is not one program: it is a
family $\textsc{Solve}_{\Gamma}$ with a polynomial $p_{\Gamma}$ for each
$\Gamma$, and \textup{(T2)} must be stated in that form. Attempting a bound
uniform in $\Gamma$ would be proving something else, and something false.
\end{hazard}

\begin{hazard}[From pseudocode to a Lean function]
\label{haz:algorithm-formalization}
\begin{enumerate}\alphlabels\tightlist
\item \emph{Termination is not structural.} $\textsc{Solve}$ recurses on
  strictly smaller domains, and $\textsc{SolveLinear}$ loops on a strictly
  decreasing dimension. Neither is a structural recursion; both need an
  explicit well-founded measure, and the two interleave ---
  $\textsc{SolveLinear}$ calls $\textsc{Solve}$ on a smaller domain, which
  calls $\textsc{SolveLinear}$ again. The measure is lexicographic:
  $\bigl(\sum_{x}|D_{x}|,\ \dim\phi\bigr)$.
\item \emph{Several steps quantify over objects that are only finitely many by
  accident.} ``If some $D_{x}$ has a strong subset $B$'' quantifies over
  subsets of $D_{x}$, over congruences on $D_{x}$, and --- through absorption
  --- over \emph{term operations}, of which there are infinitely many. See
  Formalization note~\ref{haz:t1-computable} for the route to decidability that does not
  depend on an uncited general arity bound.
\item \emph{Polynomial time is a different subject.} The bound in
  \cite{ZhukJACM} counts calls and bounds the recursion depth by $|A|$; making
  that formal needs a cost model. See \S\ref{sec:complexity-layer}.
\end{enumerate}
\end{hazard}

\section{The complexity layer}\label{sec:complexity-layer}

Targets \textup{(T2)} and \textup{(H1)} are stated here and not proved. The
point of stating them is to be precise about what is \emph{not} being claimed.

\begin{definition}[What a cost model would have to supply]\label{def:cost-model}
To state \textup{(T2)} one needs: an encoding of instances of $\CSP(\Gamma)$ as
strings; a machine model with a step count; a predicate ``$f$ is computed in
time bounded by a polynomial in the input length''; and closure of that
predicate under composition. To state \textup{(H1)} one needs, in addition, a
polynomial-time many-one reduction relation $\le_{p}$ and the class
\textup{NP}.
\end{definition}

Mathlib supplies the machine (\texttt{Turing.TM2}), an encoding
(\texttt{Computability.Encoding}), and a time-bounded computation predicate
(\texttt{Turing.TM2ComputableInPolyTime}). It does not supply a complexity
\emph{class}, \textup{P}, \textup{NP}, \textup{NP}-completeness, $\le_{p}$, or
\textsc{Sat}; and the composition lemma
\texttt{TM2ComputableInPolyTime.comp} is an open \texttt{proof\_wanted} in
Mathlib itself. Moreover \texttt{TM2ComputableInPolyTime} applies to total
functions $\alpha\to\beta$, never to languages, so even the existing predicate
would need reshaping.

\begin{remark}[This layer now exists, and it was cheap]\label{rem:complexity-scope}
An earlier draft of this remark said that building the complexity layer was
``a substantial project --- comparable in size to the algebraic core'' and that
attempting it as part of this project ``would be a scoping error''. That was
wrong, and it has been tested rather than argued: the vocabulary --- words, a
cost model, $\mathrm P$, $\mathrm{NP}$ by verifiers, $\le_{p}$,
\textup{NP}-hardness, \textup{NP}-completeness, and a generic
\textup{NP}-hard problem --- is about $850$ lines of Lean, \texttt{sorry}-free.

So the reason to keep \textup{(T2)} out of scope is \emph{not} that complexity
theory is unavailable. It is Formalization note~\ref{haz:t2-is-not-small}: \textup{(T2)}
requires exhibiting $\textsc{Solve}$ as a program with a proved step bound, and
that is the largest program-construction task in the project. \textup{(H1)},
by contrast, needs only a syntactic substitution and should be built.
\end{remark}

\begin{remark}[What the exercise cost, and what it caught]
\label{rem:complexity-defects}
Two defects surfaced in those $850$ lines, both in \emph{statements} rather
than proofs, and neither of a kind a type-checker reports. The generic problem
was first defined with a single budget bounding both the certificate length and
the running time, which makes the theorem false --- the two directions pull the
bound opposite ways. And the language was first given only the tests ``register
is empty'' and ``register's first bit is $1$'', with which no loop that
consumes its input can be written at all; that was found by trying to write the
first program.

The ratio is worth noting for the rest of this project. Two statement-level
errors in $850$ lines, against fourteen found in Zhuk's fifty pages by readers
who did not formalize. The defect list of the audit document should be expected
to grow under formalization, not to be complete.
\end{remark}

\begin{remark}[\textup{NP} membership]\label{rem:np-membership}
The proof paper's hard half concludes \textup{NP}-hardness. Recovering
\textup{NP}-completeness needs the standard theorem that every fixed
finite-template CSP lies in \textup{NP} under the chosen encoding: a solution
is a certificate of size linear in the number of variables, and checking it
means evaluating each constraint, which for fixed $\Gamma$ is a table lookup.
It is small, it needs the cost model, and it belongs with \textup{(H1)}.
\end{remark}
